
\documentclass[9pt]{article}
\usepackage[margin=0.6in]{geometry}
\usepackage{enumitem}
\usepackage{array}
\usepackage{xcolor}
\usepackage{hyperref}
\usepackage{booktabs}
\usepackage{fancyhdr}
\usepackage{titlesec}

\hypersetup{colorlinks=true, linkcolor=blue}
\pagestyle{fancy}
\fancyhf{}
\fancyhead[L]{\textbf{Question Difficulty Validation Report}}
\fancyhead[R]{\thepage}
\fancyfoot[C]{}

\titleformat{\section}{\small\bfseries}{}{0em}{}
\setlist[enumerate]{label=\alph.), leftmargin=*, itemsep=0pt, topsep=0pt, parsep=0pt}
\setlength{\parindent}{0pt}
\setlength{\parskip}{3pt}

\begin{document}
\begin{center}
{\Large \textbf{Question Difficulty Validation Report}} \\[4pt]
{\small \# of Questions: 68} \\[4pt]
\noindent\rule{\linewidth}{0.4pt}
\end{center}
\vspace{0.3cm}

\section*{Question 1}

\textbf{Q:} Why is the information gain measure considered biased when evaluating categorical attributes with many distinct values?

\begin{enumerate}[label=\textbf{\alph{*})}, nosep]
    \item A) It underestimates the predictive power of attributes with few values.
    \item \textbf{B) It favors attributes that split the data into many small groups.} % CORRECT
    \item C) It assigns equal importance to all attribute values regardless of frequency.
    \item D) It ignores the hierarchical relationships among attribute values.

\end{enumerate}

\begin{tabular}{@{}l p{10.5cm}@{}}
\textbf{Answer:} & \textbf{B} \hspace{1em} Bloom: Understand \hspace{1em} Comp: 0.638s \hspace{1em} Tokens: 592 \\
\textbf{Empirical:} & \textbf{Medium} \hspace{1em} \textbf{Actual (Human):} \underline{\hspace{5cm}}
\end{tabular}

\vspace{0.2cm}
\rule{\textwidth}{0.4pt}
\vspace{0.2cm}


\section*{Question 2}

\textbf{Q:} What mechanism does the gain ratio use to mitigate the bias of information gain toward attributes with many values?

\begin{enumerate}[label=\textbf{\alph{*})}, nosep]
    \item A) It normalizes information gain by the total number of records.
    \item B) It multiplies information gain by the number of distinct attribute values.
    \item \textbf{C) It divides information gain by the intrinsic information of the attribute.} % CORRECT
    \item D) It replaces information gain with the correlation coefficient.

\end{enumerate}

\begin{tabular}{@{}l p{10.5cm}@{}}
\textbf{Answer:} & \textbf{C} \hspace{1em} Bloom: Understand \hspace{1em} Comp: 0.647s \hspace{1em} Tokens: 506 \\
\textbf{Empirical:} & \textbf{Medium} \hspace{1em} \textbf{Actual (Human):} \underline{\hspace{5cm}}
\end{tabular}

\vspace{0.2cm}
\rule{\textwidth}{0.4pt}
\vspace{0.2cm}


\section*{Question 3}

\textbf{Q:} In attribute‑oriented induction, how does a concept hierarchy help in handling categorical attributes during generalization?

\begin{enumerate}[label=\textbf{\alph{*})}, nosep]
    \item A) It converts categorical values into numeric scores for clustering.
    \item \textbf{B) It provides a predefined ladder of abstract levels to roll up specific categories.} % CORRECT
    \item C) It merges all categories into a single “unknown” value to simplify analysis.
    \item D) It eliminates the need for any threshold settings on categorical attributes.

\end{enumerate}

\begin{tabular}{@{}l p{10.5cm}@{}}
\textbf{Answer:} & \textbf{B} \hspace{1em} Bloom: Understand \hspace{1em} Comp: 0.519s \hspace{1em} Tokens: 491 \\
\textbf{Empirical:} & \textbf{Easy} \hspace{1em} \textbf{Actual (Human):} \underline{\hspace{5cm}}
\end{tabular}

\vspace{0.2cm}
\rule{\textwidth}{0.4pt}
\vspace{0.2cm}


\section*{Question 4}

\textbf{Q:} What primary benefit do join indices provide when identifying subcubes that involve categorical dimensions in a star schema?

\begin{enumerate}[label=\textbf{\alph{*})}, nosep]
    \item A) They automatically compute aggregate measures for each dimension.
    \item B) They store pre‑aggregated numeric summaries to speed up queries.
    \item \textbf{C) They maintain direct links between dimension attribute values and fact rows, enabling fast subcube retrieval.} % CORRECT
    \item D) They eliminate the need for foreign key constraints between tables.

\end{enumerate}

\begin{tabular}{@{}l p{10.5cm}@{}}
\textbf{Answer:} & \textbf{C} \hspace{1em} Bloom: Understand \hspace{1em} Comp: 0.629s \hspace{1em} Tokens: 542 \\
\textbf{Empirical:} & \textbf{Medium} \hspace{1em} \textbf{Actual (Human):} \underline{\hspace{5cm}}
\end{tabular}

\vspace{0.2cm}
\rule{\textwidth}{0.4pt}
\vspace{0.2cm}


\section*{Question 5}

\textbf{Q:} In OLAP terminology, why is the function count() treated as a measure while other attributes (often categorical) are treated as dimensions?

\begin{enumerate}[label=\textbf{\alph{*})}, nosep]
    \item \textbf{A) Because count() always returns a numeric value that can be aggregated across dimensions.} % CORRECT
    \item B) Because count() modifies the underlying categorical attribute values.
    \item C) Because dimensions cannot be used in any aggregation functions.
    \item D) Because measures are required to be non‑numeric for categorical analysis.

\end{enumerate}

\begin{tabular}{@{}l p{10.5cm}@{}}
\textbf{Answer:} & \textbf{A} \hspace{1em} Bloom: Understand \hspace{1em} Comp: 0.296s \hspace{1em} Tokens: 282 \\
\textbf{Empirical:} & \textbf{Easy} \hspace{1em} \textbf{Actual (Human):} \underline{\hspace{5cm}}
\end{tabular}

\vspace{0.2cm}
\rule{\textwidth}{0.4pt}
\vspace{0.2cm}


\section*{Question 6}

\textbf{Q:} How does reducing an attribute’s threshold influence the generalization of a categorical attribute in attribute‑oriented induction?

\begin{enumerate}[label=\textbf{\alph{*})}, nosep]
    \item \textbf{A) It causes the algorithm to roll up the attribute to a higher (more abstract) level in its hierarchy.} % CORRECT
    \item B) It forces the attribute to be treated as a numeric measure instead of a dimension.
    \item C) It eliminates the attribute from the working relation entirely.
    \item D) It increases the number of distinct values retained, preventing any generalization.

\end{enumerate}

\begin{tabular}{@{}l p{10.5cm}@{}}
\textbf{Answer:} & \textbf{A} \hspace{1em} Bloom: Understand \hspace{1em} Comp: 0.989s \hspace{1em} Tokens: 899 \\
\textbf{Empirical:} & \textbf{Hard} \hspace{1em} \textbf{Actual (Human):} \underline{\hspace{5cm}}
\end{tabular}

\vspace{0.2cm}
\rule{\textwidth}{0.4pt}
\vspace{0.2cm}


\section*{Question 7}

\textbf{Q:} Why might different attribute relevance assessment methods yield varying rankings for categorical attributes?

\begin{enumerate}[label=\textbf{\alph{*})}, nosep]
    \item \textbf{A) Because each method uses a distinct bias, such as favoring many‑valued attributes or accounting for value probabilities.} % CORRECT
    \item B) Because categorical attributes cannot be measured statistically.
    \item C) Because all methods treat categorical attributes identically, leading to identical rankings.
    \item D) Because relevance assessment only applies to numeric attributes.

\end{enumerate}

\begin{tabular}{@{}l p{10.5cm}@{}}
\textbf{Answer:} & \textbf{A} \hspace{1em} Bloom: Understand \hspace{1em} Comp: 0.382s \hspace{1em} Tokens: 331 \\
\textbf{Empirical:} & \textbf{Easy} \hspace{1em} \textbf{Actual (Human):} \underline{\hspace{5cm}}
\end{tabular}

\vspace{0.2cm}
\rule{\textwidth}{0.4pt}
\vspace{0.2cm}


\section*{Question 8}

\textbf{Q:} When scanning a relation to collect statistics for categorical attributes, which factor mainly determines the cost of computing mapping pairs (v, v′) for each attribute?

\begin{enumerate}[label=\textbf{\alph{*})}, nosep]
    \item A) The total number of rows in the fact table.
    \item \textbf{B) The number of distinct values the categorical attribute possesses.} % CORRECT
    \item C) The depth of the attribute’s concept hierarchy.
    \item D) Whether the attribute is defined as a primary key.

\end{enumerate}

\begin{tabular}{@{}l p{10.5cm}@{}}
\textbf{Answer:} & \textbf{B} \hspace{1em} Bloom: Understand \hspace{1em} Comp: 0.678s \hspace{1em} Tokens: 613 \\
\textbf{Empirical:} & \textbf{Medium} \hspace{1em} \textbf{Actual (Human):} \underline{\hspace{5cm}}
\end{tabular}

\vspace{0.2cm}
\rule{\textwidth}{0.4pt}
\vspace{0.2cm}


\section*{Question 9}

\textbf{Q:} How does a foreign‑key relationship support the handling of categorical attributes in join indexing across multiple dimensions?

\begin{enumerate}[label=\textbf{\alph{*})}, nosep]
    \item A) It ensures that every categorical value appears in all dimension tables.
    \item \textbf{B) It provides the link between a fact table’s foreign key and a dimension table’s primary key, enabling efficient joins on categorical attributes.} % CORRECT
    \item C) It automatically aggregates categorical attributes into numeric measures.
    \item D) It eliminates the need for concept hierarchies when generalizing categorical data.

\end{enumerate}

\begin{tabular}{@{}l p{10.5cm}@{}}
\textbf{Answer:} & \textbf{B} \hspace{1em} Bloom: Understand \hspace{1em} Comp: 0.418s \hspace{1em} Tokens: 290 \\
\textbf{Empirical:} & \textbf{Easy} \hspace{1em} \textbf{Actual (Human):} \underline{\hspace{5cm}}
\end{tabular}

\vspace{0.2cm}
\rule{\textwidth}{0.4pt}
\vspace{0.2cm}


\section*{Question 10}

\textbf{Q:} In attribute‑oriented induction, what effect does setting a high attribute generalization threshold have on a categorical attribute?

\begin{enumerate}[label=\textbf{\alph{*})}, nosep]
    \item A) It forces the attribute to be generalized to the most abstract level immediately.
    \item \textbf{B) It restricts the amount of generalization, keeping the attribute at a more detailed level.} % CORRECT
    \item C) It converts the categorical attribute into a numeric one.
    \item D) It removes the attribute from consideration in the mining process.

\end{enumerate}

\begin{tabular}{@{}l p{10.5cm}@{}}
\textbf{Answer:} & \textbf{B} \hspace{1em} Bloom: Understand \hspace{1em} Comp: 0.642s \hspace{1em} Tokens: 530 \\
\textbf{Empirical:} & \textbf{Medium} \hspace{1em} \textbf{Actual (Human):} \underline{\hspace{5cm}}
\end{tabular}

\vspace{0.2cm}
\rule{\textwidth}{0.4pt}
\vspace{0.2cm}


\section*{Question 11}

\textbf{Q:} How does attribute‑oriented induction potentially limit the usefulness of accuracy as a performance metric?

\begin{enumerate}[label=\textbf{\alph{*})}, nosep]
    \item A) It requires scanning the relation multiple times, increasing error.
    \item \textbf{B) Automatic generalization may hide detailed patterns, reducing the classifier’s ability to capture subtle differences.} % CORRECT
    \item C) It forces the use of chi‑square tests for nominal data.
    \item D) It guarantees perfect accuracy by removing noisy attributes.

\end{enumerate}

\begin{tabular}{@{}l p{10.5cm}@{}}
\textbf{Answer:} & \textbf{B} \hspace{1em} Bloom: Understand \hspace{1em} Comp: 0.566s \hspace{1em} Tokens: 401 \\
\textbf{Empirical:} & \textbf{Easy} \hspace{1em} \textbf{Actual (Human):} \underline{\hspace{5cm}}
\end{tabular}

\vspace{0.2cm}
\rule{\textwidth}{0.4pt}
\vspace{0.2cm}


\section*{Question 12}

\textbf{Q:} What core assumption do clustering‑based outlier detection methods make about normal data versus outliers?

\begin{enumerate}[label=\textbf{\alph{*})}, nosep]
    \item \textbf{A) Normal data belong to large and dense clusters, while outliers belong to small or sparse clusters or to no cluster.} % CORRECT
    \item B) Normal data are uniformly distributed, and outliers form a single dense cluster.
    \item C) Normal data have higher variance than outliers.
    \item D) Normal data are always located near the centroid of the dataset.

\end{enumerate}

\begin{tabular}{@{}l p{10.5cm}@{}}
\textbf{Answer:} & \textbf{A} \hspace{1em} Bloom: Understand \hspace{1em} Comp: 0.602s \hspace{1em} Tokens: 544 \\
\textbf{Empirical:} & \textbf{Medium} \hspace{1em} \textbf{Actual (Human):} \underline{\hspace{5cm}}
\end{tabular}

\vspace{0.2cm}
\rule{\textwidth}{0.4pt}
\vspace{0.2cm}


\section*{Question 13}

\textbf{Q:} Why is determining whether a new object is an outlier typically fast after clustering has been performed?

\begin{enumerate}[label=\textbf{\alph{*})}, nosep]
    \item A) Because each object is compared only to the centroids of a few clusters rather than to all individual data points.
    \item B) Because outlier detection uses a pre‑computed lookup table of all possible outliers.
    \item \textbf{C) Because the number of clusters is usually small compared to the total number of objects, reducing the number of comparisons needed.} % CORRECT
    \item D) Because outlier detection algorithms run in constant time regardless of dataset size.

\end{enumerate}

\begin{tabular}{@{}l p{10.5cm}@{}}
\textbf{Answer:} & \textbf{C} \hspace{1em} Bloom: Understand \hspace{1em} Comp: 0.783s \hspace{1em} Tokens: 640 \\
\textbf{Empirical:} & \textbf{Hard} \hspace{1em} \textbf{Actual (Human):} \underline{\hspace{5cm}}
\end{tabular}

\vspace{0.2cm}
\rule{\textwidth}{0.4pt}
\vspace{0.2cm}


\section*{Question 14}

\textbf{Q:} How does the density measure used for outlier detection differ from the density measure used in density‑based clustering?

\begin{enumerate}[label=\textbf{\alph{*})}, nosep]
    \item \textbf{A) The outlier detection density focuses on the overall size of a cluster, while density‑based clustering considers local point density in a region.} % CORRECT
    \item B) The outlier detection density evaluates only the distance to the nearest centroid, whereas density‑based clustering uses pairwise distances.
    \item C) Both measures are identical; the text states there is no difference.
    \item D) The outlier detection density is calculated after clustering, while density‑based clustering never uses density.

\end{enumerate}

\begin{tabular}{@{}l p{10.5cm}@{}}
\textbf{Answer:} & \textbf{A} \hspace{1em} Bloom: Understand \hspace{1em} Comp: 1.143s \hspace{1em} Tokens: 1000 \\
\textbf{Empirical:} & \textbf{Hard} \hspace{1em} \textbf{Actual (Human):} \underline{\hspace{5cm}}
\end{tabular}

\vspace{0.2cm}
\rule{\textwidth}{0.4pt}
\vspace{0.2cm}


\section*{Question 15}

\textbf{Q:} What is a potential limitation of using only cluster size as a criterion for outlier detection?

\begin{enumerate}[label=\textbf{\alph{*})}, nosep]
    \item \textbf{A) It may incorrectly label rare but legitimate sub‑populations as outliers.} % CORRECT
    \item B) It requires computing the exact distance between every pair of points.
    \item C) It cannot be applied to datasets with more than three dimensions.
    \item D) It always leads to higher computational complexity than other methods.

\end{enumerate}

\begin{tabular}{@{}l p{10.5cm}@{}}
\textbf{Answer:} & \textbf{A} \hspace{1em} Bloom: Understand \hspace{1em} Comp: 0.446s \hspace{1em} Tokens: 427 \\
\textbf{Empirical:} & \textbf{Easy} \hspace{1em} \textbf{Actual (Human):} \underline{\hspace{5cm}}
\end{tabular}

\vspace{0.2cm}
\rule{\textwidth}{0.4pt}
\vspace{0.2cm}


\section*{Question 16}

\textbf{Q:} Why does the small number of clusters relative to total objects contribute to efficiency in outlier detection?

\begin{enumerate}[label=\textbf{\alph{*})}, nosep]
    \item \textbf{A) Fewer clusters mean fewer distance calculations when assigning a new object to a cluster.} % CORRECT
    \item B) A small number of clusters allows the use of hierarchical indexing structures.
    \item C) It enables parallel processing of each cluster on separate CPUs.
    \item D) It reduces the need for any similarity measures altogether.

\end{enumerate}

\begin{tabular}{@{}l p{10.5cm}@{}}
\textbf{Answer:} & \textbf{A} \hspace{1em} Bloom: Understand \hspace{1em} Comp: 0.368s \hspace{1em} Tokens: 332 \\
\textbf{Empirical:} & \textbf{Easy} \hspace{1em} \textbf{Actual (Human):} \underline{\hspace{5cm}}
\end{tabular}

\vspace{0.2cm}
\rule{\textwidth}{0.4pt}
\vspace{0.2cm}


\section*{Question 17}

\textbf{Q:} What issue arises when cluster validation relies solely on the density measure discussed for outlier detection?

\begin{enumerate}[label=\textbf{\alph{*})}, nosep]
    \item \textbf{A) The measure may confuse genuinely sparse but valid clusters with outlier groups, leading to false positives.} % CORRECT
    \item B) It forces all clusters to have the same number of points.
    \item C) It eliminates the need for any further statistical testing.
    \item D) It guarantees that all identified outliers are true anomalies.

\end{enumerate}

\begin{tabular}{@{}l p{10.5cm}@{}}
\textbf{Answer:} & \textbf{A} \hspace{1em} Bloom: Understand \hspace{1em} Comp: 0.269s \hspace{1em} Tokens: 232 \\
\textbf{Empirical:} & \textbf{Easy} \hspace{1em} \textbf{Actual (Human):} \underline{\hspace{5cm}}
\end{tabular}

\vspace{0.2cm}
\rule{\textwidth}{0.4pt}
\vspace{0.2cm}


\section*{Question 18}

\textbf{Q:} Which of the following is NOT listed as an application of data mining in the provided material?

\begin{enumerate}[label=\textbf{\alph{*})}, nosep]
    \item A) Constructing analytical models
    \item B) Performing classification and prediction
    \item \textbf{C) Generating real‑time transaction processing} % CORRECT
    \item D) Discovering hidden patterns and associations

\end{enumerate}

\begin{tabular}{@{}l p{10.5cm}@{}}
\textbf{Answer:} & \textbf{C} \hspace{1em} Bloom: Remember \hspace{1em} Comp: 0.596s \hspace{1em} Tokens: 496 \\
\textbf{Empirical:} & \textbf{Medium} \hspace{1em} \textbf{Actual (Human):} \underline{\hspace{5cm}}
\end{tabular}

\vspace{0.2cm}
\rule{\textwidth}{0.4pt}
\vspace{0.2cm}


\section*{Question 19}

\textbf{Q:} How does the context differentiate decision‑support data from data stored in operational databases?

\begin{enumerate}[label=\textbf{\alph{*})}, nosep]
    \item A) Decision‑support data is raw, unclean, and lacks historic records.
    \item \textbf{B) Decision‑support data is high‑quality, clean, integrated, and often historic, whereas operational databases contain detailed raw transactions.} % CORRECT
    \item C) Operational databases contain aggregated summaries, while decision‑support data holds individual transaction logs.
    \item D) Both contain identical data; the difference lies only in the user interface.

\end{enumerate}

\begin{tabular}{@{}l p{10.5cm}@{}}
\textbf{Answer:} & \textbf{B} \hspace{1em} Bloom: Understand \hspace{1em} Comp: 0.517s \hspace{1em} Tokens: 483 \\
\textbf{Empirical:} & \textbf{Medium} \hspace{1em} \textbf{Actual (Human):} \underline{\hspace{5cm}}
\end{tabular}

\vspace{0.2cm}
\rule{\textwidth}{0.4pt}
\vspace{0.2cm}


\section*{Question 20}

\textbf{Q:} In the three‑tier architecture described, which tier is associated with performing data mining activities?

\begin{enumerate}[label=\textbf{\alph{*})}, nosep]
    \item A) Bottom tier – Data warehouse
    \item B) Middle tier – OLAP server
    \item \textbf{C) Top tier – Front‑end tools} % CORRECT
    \item D) All tiers equally share data mining responsibilities

\end{enumerate}

\begin{tabular}{@{}l p{10.5cm}@{}}
\textbf{Answer:} & \textbf{C} \hspace{1em} Bloom: Remember \hspace{1em} Comp: 0.464s \hspace{1em} Tokens: 405 \\
\textbf{Empirical:} & \textbf{Medium} \hspace{1em} \textbf{Actual (Human):} \underline{\hspace{5cm}}
\end{tabular}

\vspace{0.2cm}
\rule{\textwidth}{0.4pt}
\vspace{0.2cm}


\section*{Question 21}

\textbf{Q:} How do data/knowledge visualization tools interact with multidimensional data mining to improve analytical outcomes?

\begin{enumerate}[label=\textbf{\alph{*})}, nosep]
    \item A) They automatically correct errors in the raw data before mining.
    \item \textbf{B) They provide facilities for drilling, pivoting, filtering, dicing, and slicing, thereby enhancing flexibility and power of mining results.} % CORRECT
    \item C) They eliminate the need for any preprocessing steps.
    \item D) They store the mined models in a separate repository without user interaction.

\end{enumerate}

\begin{tabular}{@{}l p{10.5cm}@{}}
\textbf{Answer:} & \textbf{B} \hspace{1em} Bloom: Understand \hspace{1em} Comp: 0.297s \hspace{1em} Tokens: 233 \\
\textbf{Empirical:} & \textbf{Easy} \hspace{1em} \textbf{Actual (Human):} \underline{\hspace{5cm}}
\end{tabular}

\vspace{0.2cm}
\rule{\textwidth}{0.4pt}
\vspace{0.2cm}


\section*{Question 22}

\textbf{Q:} Which application of data mining is highlighted as enhancing flexibility by allowing analysts to drill, pivot, filter, dice, and slice on a data cube?

\begin{enumerate}[label=\textbf{\alph{*})}, nosep]
    \item A) Text mining for sentiment analysis
    \item \textbf{B) Multidimensional data mining} % CORRECT
    \item C) Association rule mining in market baskets
    \item D) Clustering of spatial data

\end{enumerate}

\begin{tabular}{@{}l p{10.5cm}@{}}
\textbf{Answer:} & \textbf{B} \hspace{1em} Bloom: Remember \hspace{1em} Comp: 0.308s \hspace{1em} Tokens: 293 \\
\textbf{Empirical:} & \textbf{Easy} \hspace{1em} \textbf{Actual (Human):} \underline{\hspace{5cm}}
\end{tabular}

\vspace{0.2cm}
\rule{\textwidth}{0.4pt}
\vspace{0.2cm}


\section*{Question 23}

\textbf{Q:} How does the “refresh” function contribute to the KDD cycle in a data‑warehouse environment?

\begin{enumerate}[label=\textbf{\alph{*})}, nosep]
    \item A) It aggregates data for multidimensional analysis.
    \item \textbf{B) It updates the warehouse with new or changed data from operational sources.} % CORRECT
    \item C) It merges similar data from heterogeneous sources into a unified format.
    \item D) It visualizes the results of classification models.

\end{enumerate}

\begin{tabular}{@{}l p{10.5cm}@{}}
\textbf{Answer:} & \textbf{B} \hspace{1em} Bloom: Understand \hspace{1em} Comp: 0.539s \hspace{1em} Tokens: 460 \\
\textbf{Empirical:} & \textbf{Easy} \hspace{1em} \textbf{Actual (Human):} \underline{\hspace{5cm}}
\end{tabular}

\vspace{0.2cm}
\rule{\textwidth}{0.4pt}
\vspace{0.2cm}


\section*{Question 24}

\textbf{Q:} Which statement correctly explains why “information processing” is not considered data mining?

\begin{enumerate}[label=\textbf{\alph{*})}, nosep]
    \item A) Information processing uses sophisticated pattern detection, whereas data mining does not.
    \item \textbf{B) Information processing only retrieves directly stored or aggregate information, not hidden patterns.} % CORRECT
    \item C) Information processing is limited to visualizing data cubes.
    \item D) Information processing requires attribute‑oriented induction before any analysis.

\end{enumerate}

\begin{tabular}{@{}l p{10.5cm}@{}}
\textbf{Answer:} & \textbf{B} \hspace{1em} Bloom: Understand \hspace{1em} Comp: 0.518s \hspace{1em} Tokens: 483 \\
\textbf{Empirical:} & \textbf{Medium} \hspace{1em} \textbf{Actual (Human):} \underline{\hspace{5cm}}
\end{tabular}

\vspace{0.2cm}
\rule{\textwidth}{0.4pt}
\vspace{0.2cm}


\section*{Question 25}

\textbf{Q:} Which of the following operations is NOT listed as a way multidimensional data mining can explore a data cube?

\begin{enumerate}[label=\textbf{\alph{*})}, nosep]
    \item A) Drilling
    \item B) Pivoting
    \item \textbf{C) Clustering} % CORRECT
    \item D) Slicing

\end{enumerate}

\begin{tabular}{@{}l p{10.5cm}@{}}
\textbf{Answer:} & \textbf{C} \hspace{1em} Bloom: Remember \hspace{1em} Comp: 0.542s \hspace{1em} Tokens: 468 \\
\textbf{Empirical:} & \textbf{Medium} \hspace{1em} \textbf{Actual (Human):} \underline{\hspace{5cm}}
\end{tabular}

\vspace{0.2cm}
\rule{\textwidth}{0.4pt}
\vspace{0.2cm}


\section*{Question 26}

\textbf{Q:} Which activity is specifically mentioned as a prerequisite before performing attribute‑oriented induction?

\begin{enumerate}[label=\textbf{\alph{*})}, nosep]
    \item A) Data cleaning
    \item \textbf{B) Data focusing} % CORRECT
    \item C) Data cube computation
    \item D) Data visualization

\end{enumerate}

\begin{tabular}{@{}l p{10.5cm}@{}}
\textbf{Answer:} & \textbf{B} \hspace{1em} Bloom: Remember \hspace{1em} Comp: 0.471s \hspace{1em} Tokens: 420 \\
\textbf{Empirical:} & \textbf{Medium} \hspace{1em} \textbf{Actual (Human):} \underline{\hspace{5cm}}
\end{tabular}

\vspace{0.2cm}
\rule{\textwidth}{0.4pt}
\vspace{0.2cm}


\section*{Question 27}

\textbf{Q:} Which statement best describes the role of OLAP in multidimensional data mining as presented?

\begin{enumerate}[label=\textbf{\alph{*})}, nosep]
    \item A) OLAP provides only storage for raw transaction logs.
    \item \textbf{B) OLAP enables dynamic selection and swapping of data mining functions on a data cube.} % CORRECT
    \item C) OLAP replaces the need for any data mining algorithms.
    \item D) OLAP limits analysis to a single predefined query.

\end{enumerate}

\begin{tabular}{@{}l p{10.5cm}@{}}
\textbf{Answer:} & \textbf{B} \hspace{1em} Bloom: Understand \hspace{1em} Comp: 0.618s \hspace{1em} Tokens: 566 \\
\textbf{Empirical:} & \textbf{Medium} \hspace{1em} \textbf{Actual (Human):} \underline{\hspace{5cm}}
\end{tabular}

\vspace{0.2cm}
\rule{\textwidth}{0.4pt}
\vspace{0.2cm}


\section*{Question 28}

\textbf{Q:} Which of the following best captures the overall purpose of data mining within the knowledge discovery process?

\begin{enumerate}[label=\textbf{\alph{*})}, nosep]
    \item A) To store operational data for daily transactions.
    \item B) To generate aggregate reports using simple SQL queries.
    \item \textbf{C) To uncover hidden patterns and construct models that transform data into actionable knowledge.} % CORRECT
    \item D) To visualize existing information without altering it.

\end{enumerate}

\begin{tabular}{@{}l p{10.5cm}@{}}
\textbf{Answer:} & \textbf{C} \hspace{1em} Bloom: Understand \hspace{1em} Comp: 0.576s \hspace{1em} Tokens: 455 \\
\textbf{Empirical:} & \textbf{Easy} \hspace{1em} \textbf{Actual (Human):} \underline{\hspace{5cm}}
\end{tabular}

\vspace{0.2cm}
\rule{\textwidth}{0.4pt}
\vspace{0.2cm}


\section*{Question 29}

\textbf{Q:} Which statement best describes a key difference between data mining and data warehousing in terms of data granularity?

\begin{enumerate}[label=\textbf{\alph{*})}, nosep]
    \item A) Data mining always works on highly summarized data, while data warehousing stores raw transaction details.
    \item \textbf{B) Data mining may analyze data at more detailed granularities than the summarized data stored in a data warehouse.} % CORRECT
    \item C) Data warehousing performs pattern discovery, whereas data mining only stores data.
    \item D) Data mining and data warehousing always use the same level of data detail.

\end{enumerate}

\begin{tabular}{@{}l p{10.5cm}@{}}
\textbf{Answer:} & \textbf{B} \hspace{1em} Bloom: Understand \hspace{1em} Comp: 0.706s \hspace{1em} Tokens: 654 \\
\textbf{Empirical:} & \textbf{Medium} \hspace{1em} \textbf{Actual (Human):} \underline{\hspace{5cm}}
\end{tabular}

\vspace{0.2cm}
\rule{\textwidth}{0.4pt}
\vspace{0.2cm}


\section*{Question 30}

\textbf{Q:} Which of the following is NOT listed as a back‑end tool or utility for populating and refreshing a data warehouse?

\begin{enumerate}[label=\textbf{\alph{*})}, nosep]
    \item A) Data extraction
    \item B) Data cleaning
    \item \textbf{C) Data indexing} % CORRECT
    \item D) Data transformation

\end{enumerate}

\begin{tabular}{@{}l p{10.5cm}@{}}
\textbf{Answer:} & \textbf{C} \hspace{1em} Bloom: Understand \hspace{1em} Comp: 0.601s \hspace{1em} Tokens: 499 \\
\textbf{Empirical:} & \textbf{Medium} \hspace{1em} \textbf{Actual (Human):} \underline{\hspace{5cm}}
\end{tabular}

\vspace{0.2cm}
\rule{\textwidth}{0.4pt}
\vspace{0.2cm}


\section*{Question 31}

\textbf{Q:} Which item would you expect to find in a data warehouse metadata repository?

\begin{enumerate}[label=\textbf{\alph{*})}, nosep]
    \item A) Detailed transaction logs of the source systems
    \item \textbf{B) Algorithms used for data summarization} % CORRECT
    \item C) User passwords for warehouse access
    \item D) Real‑time sensor data streams

\end{enumerate}

\begin{tabular}{@{}l p{10.5cm}@{}}
\textbf{Answer:} & \textbf{B} \hspace{1em} Bloom: Understand \hspace{1em} Comp: 1.079s \hspace{1em} Tokens: 1000 \\
\textbf{Empirical:} & \textbf{Hard} \hspace{1em} \textbf{Actual (Human):} \underline{\hspace{5cm}}
\end{tabular}

\vspace{0.2cm}
\rule{\textwidth}{0.4pt}
\vspace{0.2cm}


\section*{Question 32}

\textbf{Q:} Considering the metadata repository of a data warehouse, which of the following statements accurately reflects its purpose?

\begin{enumerate}[label=\textbf{\alph{*})}, nosep]
    \item A) It stores only the raw source data before transformation.
    \item \textbf{B) It provides details about warehouse structure, data history, summarization algorithms, and source‑to‑warehouse mappings.} % CORRECT
    \item C) It records user login attempts and security breaches.
    \item D) It contains the actual business transaction records.

\end{enumerate}

\begin{tabular}{@{}l p{10.5cm}@{}}
\textbf{Answer:} & \textbf{B} \hspace{1em} Bloom: Understand \hspace{1em} Comp: 0.436s \hspace{1em} Tokens: 409 \\
\textbf{Empirical:} & \textbf{Easy} \hspace{1em} \textbf{Actual (Human):} \underline{\hspace{5cm}}
\end{tabular}

\vspace{0.2cm}
\rule{\textwidth}{0.4pt}
\vspace{0.2cm}


\section*{Question 33}

\textbf{Q:} Which statement correctly distinguishes the granularity of data that data mining can handle compared to that stored in a data warehouse?

\begin{enumerate}[label=\textbf{\alph{*})}, nosep]
    \item A) Data mining only works with highly summarized data, while a data warehouse stores detailed transaction records.
    \item \textbf{B) Data mining can analyze both detailed and summarized data, whereas a data warehouse typically stores summarized data for decision support.} % CORRECT
    \item C) Data mining stores raw data, and a data warehouse performs the mining analysis.
    \item D) Both data mining and data warehouses always operate on the same level of data detail.

\end{enumerate}

\begin{tabular}{@{}l p{10.5cm}@{}}
\textbf{Answer:} & \textbf{B} \hspace{1em} Bloom: Understand \hspace{1em} Comp: 0.703s \hspace{1em} Tokens: 628 \\
\textbf{Empirical:} & \textbf{Medium} \hspace{1em} \textbf{Actual (Human):} \underline{\hspace{5cm}}
\end{tabular}

\vspace{0.2cm}
\rule{\textwidth}{0.4pt}
\vspace{0.2cm}


\section*{Question 34}

\textbf{Q:} Why are data warehouses considered valuable sources for data‑mining tools?

\begin{enumerate}[label=\textbf{\alph{*})}, nosep]
    \item A) They generate mining algorithms automatically.
    \item \textbf{B) They provide integrated, consistent, and cleaned data after extensive preprocessing.} % CORRECT
    \item C) They contain only raw transaction logs without any transformation.
    \item D) They restrict access to only summary statistics, preventing detailed analysis.

\end{enumerate}

\begin{tabular}{@{}l p{10.5cm}@{}}
\textbf{Answer:} & \textbf{B} \hspace{1em} Bloom: Analyze \hspace{1em} Comp: 0.449s \hspace{1em} Tokens: 418 \\
\textbf{Empirical:} & \textbf{Easy} \hspace{1em} \textbf{Actual (Human):} \underline{\hspace{5cm}}
\end{tabular}

\vspace{0.2cm}
\rule{\textwidth}{0.4pt}
\vspace{0.2cm}


\section*{Question 35}

\textbf{Q:} Which characteristic differentiates a data warehouse from a traditional RDBMS/database?

\begin{enumerate}[label=\textbf{\alph{*})}, nosep]
    \item A) Data warehouses are transaction‑oriented and frequently updated.
    \item \textbf{B) Data warehouses are subject‑oriented, integrated, time‑variant, and non‑volatile.} % CORRECT
    \item C) RDBMS stores only summarized data for decision making.
    \item D) Data warehouses do not support any form of querying.

\end{enumerate}

\begin{tabular}{@{}l p{10.5cm}@{}}
\textbf{Answer:} & \textbf{B} \hspace{1em} Bloom: Understand \hspace{1em} Comp: 0.567s \hspace{1em} Tokens: 495 \\
\textbf{Empirical:} & \textbf{Medium} \hspace{1em} \textbf{Actual (Human):} \underline{\hspace{5cm}}
\end{tabular}

\vspace{0.2cm}
\rule{\textwidth}{0.4pt}
\vspace{0.2cm}


\section*{Question 36}

\textbf{Q:} How does the metadata repository of a data warehouse support data‑mining activities?

\begin{enumerate}[label=\textbf{\alph{*})}, nosep]
    \item A) It stores the raw source files used for mining.
    \item \textbf{B) It provides details about warehouse structure, data history, summarization algorithms, and source‑to‑warehouse mappings.} % CORRECT
    \item C) It only tracks user login information.
    \item D) It replaces the need for any data preprocessing before mining.

\end{enumerate}

\begin{tabular}{@{}l p{10.5cm}@{}}
\textbf{Answer:} & \textbf{B} \hspace{1em} Bloom: Analyze \hspace{1em} Comp: 0.218s \hspace{1em} Tokens: 191 \\
\textbf{Empirical:} & \textbf{Easy} \hspace{1em} \textbf{Actual (Human):} \underline{\hspace{5cm}}
\end{tabular}

\vspace{0.2cm}
\rule{\textwidth}{0.4pt}
\vspace{0.2cm}


\section*{Question 37}

\textbf{Q:} Which of the following data types is explicitly mentioned as being outside the typical scope of a data warehouse but within the scope of data mining?

\begin{enumerate}[label=\textbf{\alph{*})}, nosep]
    \item A) Summarized sales totals per country.
    \item \textbf{B) Flat files containing unstructured text.} % CORRECT
    \item C) Relational tables with foreign keys.
    \item D) Time‑variant summary tables.

\end{enumerate}

\begin{tabular}{@{}l p{10.5cm}@{}}
\textbf{Answer:} & \textbf{B} \hspace{1em} Bloom: Understand \hspace{1em} Comp: 0.598s \hspace{1em} Tokens: 534 \\
\textbf{Empirical:} & \textbf{Easy} \hspace{1em} \textbf{Actual (Human):} \underline{\hspace{5cm}}
\end{tabular}

\vspace{0.2cm}
\rule{\textwidth}{0.4pt}
\vspace{0.2cm}


\section*{Question 38}

\textbf{Q:} In what way do ETL (Extraction, Transformation, Loading) processes influence the readiness of data for mining in a data warehouse?

\begin{enumerate}[label=\textbf{\alph{*})}, nosep]
    \item A) ETL processes delay mining by storing only raw data.
    \item \textbf{B) They clean, integrate, and transform data, producing high‑quality, consistent datasets suitable for mining.} % CORRECT
    \item C) ETL automatically discovers hidden patterns without further analysis.
    \item D) They convert mined patterns back into operational databases.

\end{enumerate}

\begin{tabular}{@{}l p{10.5cm}@{}}
\textbf{Answer:} & \textbf{B} \hspace{1em} Bloom: Analyze \hspace{1em} Comp: 0.725s \hspace{1em} Tokens: 542 \\
\textbf{Empirical:} & \textbf{Medium} \hspace{1em} \textbf{Actual (Human):} \underline{\hspace{5cm}}
\end{tabular}

\vspace{0.2cm}
\rule{\textwidth}{0.4pt}
\vspace{0.2cm}


\section*{Question 39}

\textbf{Q:} Which of the following best explains why data normalization is applied before distance‑based mining algorithms?

\begin{enumerate}[label=\textbf{\alph{*})}, nosep]
    \item A) It reduces the number of attributes in the dataset.
    \item \textbf{B) It scales all attribute values to a common range, preventing attributes with larger scales from dominating distance calculations.} % CORRECT
    \item C) It converts categorical attributes into numeric ones.
    \item D) It guarantees that all values become integers.

\end{enumerate}

\begin{tabular}{@{}l p{10.5cm}@{}}
\textbf{Answer:} & \textbf{B} \hspace{1em} Bloom: Understand \hspace{1em} Comp: 0.550s \hspace{1em} Tokens: 501 \\
\textbf{Empirical:} & \textbf{Medium} \hspace{1em} \textbf{Actual (Human):} \underline{\hspace{5cm}}
\end{tabular}

\vspace{0.2cm}
\rule{\textwidth}{0.4pt}
\vspace{0.2cm}


\section*{Question 40}

\textbf{Q:} After applying min‑max normalization, what is the possible value range for any attribute?

\begin{enumerate}[label=\textbf{\alph{*})}, nosep]
    \item A) (−∞, +∞)
    \item \textbf{B) (0, 1) inclusive} % CORRECT
    \item C) (−1, 1) inclusive
    \item D) (0, 10) inclusive

\end{enumerate}

\begin{tabular}{@{}l p{10.5cm}@{}}
\textbf{Answer:} & \textbf{B} \hspace{1em} Bloom: Understand \hspace{1em} Comp: 0.596s \hspace{1em} Tokens: 563 \\
\textbf{Empirical:} & \textbf{Medium} \hspace{1em} \textbf{Actual (Human):} \underline{\hspace{5cm}}
\end{tabular}

\vspace{0.2cm}
\rule{\textwidth}{0.4pt}
\vspace{0.2cm}


\section*{Question 41}

\textbf{Q:} Which interval correctly describes the theoretical range of values produced by z‑score normalization?

\begin{enumerate}[label=\textbf{\alph{*})}, nosep]
    \item A) (0, 1)
    \item B) (−1, 1)
    \item \textbf{C) (−∞, +∞)} % CORRECT
    \item D) (0, 10)

\end{enumerate}

\begin{tabular}{@{}l p{10.5cm}@{}}
\textbf{Answer:} & \textbf{C} \hspace{1em} Bloom: Understand \hspace{1em} Comp: 0.613s \hspace{1em} Tokens: 570 \\
\textbf{Empirical:} & \textbf{Medium} \hspace{1em} \textbf{Actual (Human):} \underline{\hspace{5cm}}
\end{tabular}

\vspace{0.2cm}
\rule{\textwidth}{0.4pt}
\vspace{0.2cm}


\section*{Question 42}

\textbf{Q:} Which of the following best describes the primary motivation for performing data preprocessing as presented in the overview?

\begin{enumerate}[label=\textbf{\alph{*})}, nosep]
    \item A) To increase the size of the dataset for better mining results
    \item \textbf{B) To improve data quality by addressing issues such as noise and inconsistency} % CORRECT
    \item C) To replace relational databases with data warehouses
    \item D) To eliminate the need for any further data analysis steps

\end{enumerate}

\begin{tabular}{@{}l p{10.5cm}@{}}
\textbf{Answer:} & \textbf{B} \hspace{1em} Bloom: Understand \hspace{1em} Comp: 0.422s \hspace{1em} Tokens: 382 \\
\textbf{Empirical:} & \textbf{Easy} \hspace{1em} \textbf{Actual (Human):} \underline{\hspace{5cm}}
\end{tabular}

\vspace{0.2cm}
\rule{\textwidth}{0.4pt}
\vspace{0.2cm}


\section*{Question 43}

\textbf{Q:} Which statement best captures the combined role of data sampling and dimensionality reduction in data reduction strategies?

\begin{enumerate}[label=\textbf{\alph{*})}, nosep]
    \item A) Both techniques create a completely new dataset that discards the original information.
    \item \textbf{B) Both aim to produce a smaller representation of the data while preserving the analytical usefulness of the original set.} % CORRECT
    \item C) Sampling reduces attributes, whereas dimensionality reduction reduces the number of tuples.
    \item D) Neither technique affects the size of the data; they only improve visualisation.

\end{enumerate}

\begin{tabular}{@{}l p{10.5cm}@{}}
\textbf{Answer:} & \textbf{B} \hspace{1em} Bloom: Understand \hspace{1em} Comp: 0.617s \hspace{1em} Tokens: 500 \\
\textbf{Empirical:} & \textbf{Medium} \hspace{1em} \textbf{Actual (Human):} \underline{\hspace{5cm}}
\end{tabular}

\vspace{0.2cm}
\rule{\textwidth}{0.4pt}
\vspace{0.2cm}


\section*{Question 44}

\textbf{Q:} Which of the following best describes the difference between parametric and non‑parametric models used in numerosity reduction?

\begin{enumerate}[label=\textbf{\alph{*})}, nosep]
    \item A) Parametric models require no assumptions about data distribution, while non‑parametric models do.
    \item \textbf{B) Parametric models assume a specific functional form (e.g., regression), whereas non‑parametric models do not (e.g., histograms, clustering).} % CORRECT
    \item C) Non‑parametric models always produce lossless compression, while parametric models are lossy.
    \item D) Parametric models are only used for dimensionality reduction, not numerosity reduction.

\end{enumerate}

\begin{tabular}{@{}l p{10.5cm}@{}}
\textbf{Answer:} & \textbf{B} \hspace{1em} Bloom: Understand \hspace{1em} Comp: 0.450s \hspace{1em} Tokens: 429 \\
\textbf{Empirical:} & \textbf{Easy} \hspace{1em} \textbf{Actual (Human):} \underline{\hspace{5cm}}
\end{tabular}

\vspace{0.2cm}
\rule{\textwidth}{0.4pt}
\vspace{0.2cm}


\section*{Question 45}

\textbf{Q:} If an attribute’s domain contains n distinct values, how many bits are required for each entry in its bitmap index?

\begin{enumerate}[label=\textbf{\alph{*})}, nosep]
    \item A) 1 bit
    \item B) log₂ n bits
    \item \textbf{C) n bits} % CORRECT
    \item D) n² bits

\end{enumerate}

\begin{tabular}{@{}l p{10.5cm}@{}}
\textbf{Answer:} & \textbf{C} \hspace{1em} Bloom: Remember \hspace{1em} Comp: 0.498s \hspace{1em} Tokens: 431 \\
\textbf{Empirical:} & \textbf{Medium} \hspace{1em} \textbf{Actual (Human):} \underline{\hspace{5cm}}
\end{tabular}

\vspace{0.2cm}
\rule{\textwidth}{0.4pt}
\vspace{0.2cm}


\section*{Question 46}

\textbf{Q:} Which of the following statements best defines a data warehouse?

\begin{enumerate}[label=\textbf{\alph{*})}, nosep]
    \item A) A volatile collection of operational data used for day‑to‑day transactions.
    \item \textbf{B) A subject‑oriented, integrated, time‑variant, and non‑volatile data collection supporting management decision making.} % CORRECT
    \item C) A repository that stores only raw, unprocessed data from source systems.
    \item D) A tool that only provides front‑end query and reporting capabilities.

\end{enumerate}

\begin{tabular}{@{}l p{10.5cm}@{}}
\textbf{Answer:} & \textbf{B} \hspace{1em} Bloom: Remember \hspace{1em} Comp: 0.695s \hspace{1em} Tokens: 534 \\
\textbf{Empirical:} & \textbf{Medium} \hspace{1em} \textbf{Actual (Human):} \underline{\hspace{5cm}}
\end{tabular}

\vspace{0.2cm}
\rule{\textwidth}{0.4pt}
\vspace{0.2cm}


\section*{Question 47}

\textbf{Q:} Which of the following is NOT listed as a back‑end tool or utility used for populating and refreshing a data warehouse?

\begin{enumerate}[label=\textbf{\alph{*})}, nosep]
    \item A) Data extraction
    \item B) Data cleaning
    \item C) Data transformation
    \item \textbf{D) Data visualization} % CORRECT

\end{enumerate}

\begin{tabular}{@{}l p{10.5cm}@{}}
\textbf{Answer:} & \textbf{D} \hspace{1em} Bloom: Remember \hspace{1em} Comp: 0.444s \hspace{1em} Tokens: 421 \\
\textbf{Empirical:} & \textbf{Medium} \hspace{1em} \textbf{Actual (Human):} \underline{\hspace{5cm}}
\end{tabular}

\vspace{0.2cm}
\rule{\textwidth}{0.4pt}
\vspace{0.2cm}


\section*{Question 48}

\textbf{Q:} Which of the following is **not** considered one of the four key characteristics of a data warehouse?

\begin{enumerate}[label=\textbf{\alph{*})}, nosep]
    \item A) Subject‑oriented
    \item B) Integrated
    \item C) Time‑variant
    \item \textbf{D) Volatile} % CORRECT

\end{enumerate}

\begin{tabular}{@{}l p{10.5cm}@{}}
\textbf{Answer:} & \textbf{D} \hspace{1em} Bloom: Remember \hspace{1em} Comp: 0.488s \hspace{1em} Tokens: 426 \\
\textbf{Empirical:} & \textbf{Easy} \hspace{1em} \textbf{Actual (Human):} \underline{\hspace{5cm}}
\end{tabular}

\vspace{0.2cm}
\rule{\textwidth}{0.4pt}
\vspace{0.2cm}


\section*{Question 49}

\textbf{Q:} Why is normalizing data to a range such as [0.0, 1.0] important when applying a distance‑based anomaly detection algorithm?

\begin{enumerate}[label=\textbf{\alph{*})}, nosep]
    \item A) It reduces the dimensionality of the data, making clustering faster.
    \item \textbf{B) It ensures that all attributes contribute equally to distance calculations, improving accuracy.} % CORRECT
    \item C) It removes outliers before the algorithm runs, guaranteeing no false alarms.
    \item D) It converts categorical attributes into numerical values suitable for distance metrics.

\end{enumerate}

\begin{tabular}{@{}l p{10.5cm}@{}}
\textbf{Answer:} & \textbf{B} \hspace{1em} Bloom: Understand \hspace{1em} Comp: 0.707s \hspace{1em} Tokens: 516 \\
\textbf{Empirical:} & \textbf{Medium} \hspace{1em} \textbf{Actual (Human):} \underline{\hspace{5cm}}
\end{tabular}

\vspace{0.2cm}
\rule{\textwidth}{0.4pt}
\vspace{0.2cm}


\section*{Question 50}

\textbf{Q:} In a clustering‑based distance outlier detection method, what indicates that an object is likely an outlier?

\begin{enumerate}[label=\textbf{\alph{*})}, nosep]
    \item A) The object lies exactly at the centroid of its nearest cluster.
    \item \textbf{B) The distance between the object and the nearest cluster’s centroid is larger than a predefined threshold.} % CORRECT
    \item C) The object belongs to multiple clusters simultaneously.
    \item D) The object has the smallest average distance to all other objects in the dataset.

\end{enumerate}

\begin{tabular}{@{}l p{10.5cm}@{}}
\textbf{Answer:} & \textbf{B} \hspace{1em} Bloom: Understand \hspace{1em} Comp: 0.616s \hspace{1em} Tokens: 584 \\
\textbf{Empirical:} & \textbf{Medium} \hspace{1em} \textbf{Actual (Human):} \underline{\hspace{5cm}}
\end{tabular}

\vspace{0.2cm}
\rule{\textwidth}{0.4pt}
\vspace{0.2cm}


\section*{Question 51}

\textbf{Q:} Which of the following best describes the role of “distance to the closest cluster” in clustering‑based outlier detection?

\begin{enumerate}[label=\textbf{\alph{*})}, nosep]
    \item A) It is used to compute the silhouette score for cluster validation.
    \item \textbf{B) It serves as a direct measure to decide whether an object deviates significantly from typical cluster patterns.} % CORRECT
    \item C) It determines the number of clusters to be formed in the dataset.
    \item D) It replaces the need for any data preprocessing such as normalization.

\end{enumerate}

\begin{tabular}{@{}l p{10.5cm}@{}}
\textbf{Answer:} & \textbf{B} \hspace{1em} Bloom: Understand \hspace{1em} Comp: 0.429s \hspace{1em} Tokens: 408 \\
\textbf{Empirical:} & \textbf{Easy} \hspace{1em} \textbf{Actual (Human):} \underline{\hspace{5cm}}
\end{tabular}

\vspace{0.2cm}
\rule{\textwidth}{0.4pt}
\vspace{0.2cm}


\section*{Question 52}

\textbf{Q:} When clustering is used to generate a concept hierarchy, each cluster becomes what element of the hierarchy?

\begin{enumerate}[label=\textbf{\alph{*})}, nosep]
    \item A) A leaf node only
    \item \textbf{B) A node of the concept hierarchy} % CORRECT
    \item C) The root of the hierarchy exclusively
    \item D) An isolated sub‑hierarchy

\end{enumerate}

\begin{tabular}{@{}l p{10.5cm}@{}}
\textbf{Answer:} & \textbf{B} \hspace{1em} Bloom: Understand \hspace{1em} Comp: 0.564s \hspace{1em} Tokens: 490 \\
\textbf{Empirical:} & \textbf{Medium} \hspace{1em} \textbf{Actual (Human):} \underline{\hspace{5cm}}
\end{tabular}

\vspace{0.2cm}
\rule{\textwidth}{0.4pt}
\vspace{0.2cm}


\section*{Question 53}

\textbf{Q:} Which approach starts with each data point as its own cluster and iteratively merges them?

\begin{enumerate}[label=\textbf{\alph{*})}, nosep]
    \item A) Top‑down splitting
    \item \textbf{B) Bottom‑up merging} % CORRECT
    \item C) K‑means clustering
    \item D) Divisive hierarchical clustering

\end{enumerate}

\begin{tabular}{@{}l p{10.5cm}@{}}
\textbf{Answer:} & \textbf{B} \hspace{1em} Bloom: Understand \hspace{1em} Comp: 0.579s \hspace{1em} Tokens: 549 \\
\textbf{Empirical:} & \textbf{Easy} \hspace{1em} \textbf{Actual (Human):} \underline{\hspace{5cm}}
\end{tabular}

\vspace{0.2cm}
\rule{\textwidth}{0.4pt}
\vspace{0.2cm}


\section*{Question 54}

\textbf{Q:} When data are “contaminated, with some values out of range or of a different data type than expected,” which phase of the data preprocessing pipeline is most appropriate to address this issue?

\begin{enumerate}[label=\textbf{\alph{*})}, nosep]
    \item A) Data integration
    \item \textbf{B) Data cleaning} % CORRECT
    \item C) Data transformation
    \item D) Data reduction

\end{enumerate}

\begin{tabular}{@{}l p{10.5cm}@{}}
\textbf{Answer:} & \textbf{B} \hspace{1em} Bloom: Understand \hspace{1em} Comp: 0.528s \hspace{1em} Tokens: 506 \\
\textbf{Empirical:} & \textbf{Medium} \hspace{1em} \textbf{Actual (Human):} \underline{\hspace{5cm}}
\end{tabular}

\vspace{0.2cm}
\rule{\textwidth}{0.4pt}
\vspace{0.2cm}


\section*{Question 55}

\textbf{Q:} Which of the following best explains why “ignoring the tuple” is a poor strategy when the missing values are limited to a few attributes rather than many?

\begin{enumerate}[label=\textbf{\alph{*})}, nosep]
    \item \textbf{A) It reduces the dataset size unnecessarily, losing valuable information.} % CORRECT
    \item B) It violates the assumption of independence among attributes.
    \item C) It increases computational complexity of subsequent mining tasks.
    \item D) It introduces duplicate records into the dataset.

\end{enumerate}

\begin{tabular}{@{}l p{10.5cm}@{}}
\textbf{Answer:} & \textbf{A} \hspace{1em} Bloom: Understand \hspace{1em} Comp: 0.470s \hspace{1em} Tokens: 449 \\
\textbf{Empirical:} & \textbf{Medium} \hspace{1em} \textbf{Actual (Human):} \underline{\hspace{5cm}}
\end{tabular}

\vspace{0.2cm}
\rule{\textwidth}{0.4pt}
\vspace{0.2cm}


\section*{Question 56}

\textbf{Q:} In a multidimensional data mining scenario, which of the following best describes the role of OLAP?

\begin{enumerate}[label=\textbf{\alph{*})}, nosep]
    \item A) It replaces the need for any data mining algorithms.
    \item \textbf{B) It provides the exploratory interface that enables mining of knowledge from multidimensional databases.} % CORRECT
    \item C) It stores only transactional (OLTP) data for fast updates.
    \item D) It limits the analysis to two‑dimensional tabular views.

\end{enumerate}

\begin{tabular}{@{}l p{10.5cm}@{}}
\textbf{Answer:} & \textbf{B} \hspace{1em} Bloom: Understand \hspace{1em} Comp: 0.664s \hspace{1em} Tokens: 560 \\
\textbf{Empirical:} & \textbf{Medium} \hspace{1em} \textbf{Actual (Human):} \underline{\hspace{5cm}}
\end{tabular}

\vspace{0.2cm}
\rule{\textwidth}{0.4pt}
\vspace{0.2cm}


\section*{Question 57}

\textbf{Q:} In the three‑tier data warehousing architecture, which tier is primarily responsible for executing OLAP operations such as roll‑up and drill‑down?

\begin{enumerate}[label=\textbf{\alph{*})}, nosep]
    \item A) The presentation tier (client).
    \item \textbf{B) The middle tier (OLAP server).} % CORRECT
    \item C) The bottom tier (data storage).
    \item D) The integration tier (ETL processes).

\end{enumerate}

\begin{tabular}{@{}l p{10.5cm}@{}}
\textbf{Answer:} & \textbf{B} \hspace{1em} Bloom: Understand \hspace{1em} Comp: 0.584s \hspace{1em} Tokens: 544 \\
\textbf{Empirical:} & \textbf{Medium} \hspace{1em} \textbf{Actual (Human):} \underline{\hspace{5cm}}
\end{tabular}

\vspace{0.2cm}
\rule{\textwidth}{0.4pt}
\vspace{0.2cm}


\section*{Question 58}

\textbf{Q:} Which OLAP operation aggregates data by moving from a lower‑level member to a higher‑level member in a dimension hierarchy (e.g., from city to province/state)?

\begin{enumerate}[label=\textbf{\alph{*})}, nosep]
    \item A) Drill‑down
    \item \textbf{B) Roll‑up} % CORRECT
    \item C) Slice
    \item D) Pivot

\end{enumerate}

\begin{tabular}{@{}l p{10.5cm}@{}}
\textbf{Answer:} & \textbf{B} \hspace{1em} Bloom: Understand \hspace{1em} Comp: 0.511s \hspace{1em} Tokens: 420 \\
\textbf{Empirical:} & \textbf{Easy} \hspace{1em} \textbf{Actual (Human):} \underline{\hspace{5cm}}
\end{tabular}

\vspace{0.2cm}
\rule{\textwidth}{0.4pt}
\vspace{0.2cm}


\section*{Question 59}

\textbf{Q:} A query requests sales for the year 2010 across the dimensions \{brand, province/state\}. Which OLAP operation(s) should be applied first to restrict the data to the year 2010?

\begin{enumerate}[label=\textbf{\alph{*})}, nosep]
    \item A) Drill‑down on year, then slice on brand
    \item \textbf{B) Slice on year, then projection on brand and province/state} % CORRECT
    \item C) Roll‑up on year, then dice on brand and province/state
    \item D) Pivot on year, then roll‑up on brand

\end{enumerate}

\begin{tabular}{@{}l p{10.5cm}@{}}
\textbf{Answer:} & \textbf{B} \hspace{1em} Bloom: Apply \hspace{1em} Comp: 0.772s \hspace{1em} Tokens: 707 \\
\textbf{Empirical:} & \textbf{Hard} \hspace{1em} \textbf{Actual (Human):} \underline{\hspace{5cm}}
\end{tabular}

\vspace{0.2cm}
\rule{\textwidth}{0.4pt}
\vspace{0.2cm}


\section*{Question 60}

\textbf{Q:} Considering the three‑tier model, which layer is most likely to be refreshed periodically through ETL processes, and how does this differ from the refresh behavior of OLTP systems?

\begin{enumerate}[label=\textbf{\alph{*})}, nosep]
    \item A) Top tier – refreshed continuously like OLTP.
    \item B) Middle tier – refreshed only on demand, unlike OLTP’s real‑time updates.
    \item C) Bottom tier – operational databases are refreshed via ETL, similar to OLTP.
    \item \textbf{D) Bottom tier – data warehouse is refreshed periodically, whereas OLTP updates are immediate.} % CORRECT

\end{enumerate}

\begin{tabular}{@{}l p{10.5cm}@{}}
\textbf{Answer:} & \textbf{D} \hspace{1em} Bloom: Understand \hspace{1em} Comp: 0.744s \hspace{1em} Tokens: 691 \\
\textbf{Empirical:} & \textbf{Hard} \hspace{1em} \textbf{Actual (Human):} \underline{\hspace{5cm}}
\end{tabular}

\vspace{0.2cm}
\rule{\textwidth}{0.4pt}
\vspace{0.2cm}


\section*{Question 61}

\textbf{Q:} Considering the data update frequency, which statement accurately reflects the difference between OLTP and OLAP environments?

\begin{enumerate}[label=\textbf{\alph{*})}, nosep]
    \item A) OLTP databases are refreshed nightly, while OLAP warehouses are updated in real time.
    \item \textbf{B) OLTP systems handle continuous, real‑time updates; OLAP systems are refreshed periodically (e.g., nightly or weekly).} % CORRECT
    \item C) Both OLTP and OLAP systems update data continuously at the same rate.
    \item D) OLAP systems never receive new data after the initial load.

\end{enumerate}

\begin{tabular}{@{}l p{10.5cm}@{}}
\textbf{Answer:} & \textbf{B} \hspace{1em} Bloom: Understand \hspace{1em} Comp: 0.522s \hspace{1em} Tokens: 433 \\
\textbf{Empirical:} & \textbf{Medium} \hspace{1em} \textbf{Actual (Human):} \underline{\hspace{5cm}}
\end{tabular}

\vspace{0.2cm}
\rule{\textwidth}{0.4pt}
\vspace{0.2cm}


\section*{Question 62}

\textbf{Q:} Which of the following best explains why OLAP systems typically require a different schema design (e.g., star or snowflake) compared to OLTP systems?

\begin{enumerate}[label=\textbf{\alph{*})}, nosep]
    \item A) OLAP schemas minimize redundancy to speed up transaction processing.
    \item B) OLTP schemas are optimized for complex joins across many dimensions.
    \item \textbf{C) OLAP schemas organize data around facts and dimensions to facilitate fast aggregation and slicing, unlike normalized OLTP schemas.} % CORRECT
    \item D) Both OLTP and OLAP use identical normalized schemas, so there is no difference.

\end{enumerate}

\begin{tabular}{@{}l p{10.5cm}@{}}
\textbf{Answer:} & \textbf{C} \hspace{1em} Bloom: Analyze \hspace{1em} Comp: 0.674s \hspace{1em} Tokens: 559 \\
\textbf{Empirical:} & \textbf{Medium} \hspace{1em} \textbf{Actual (Human):} \underline{\hspace{5cm}}
\end{tabular}

\vspace{0.2cm}
\rule{\textwidth}{0.4pt}
\vspace{0.2cm}


\section*{Question 63}

\textbf{Q:} Which of the following best explains why the choice of a deviation measurement is critical for detecting global outliers?

\begin{enumerate}[label=\textbf{\alph{*})}, nosep]
    \item A) Because the measurement determines the dimensionality of the data representation
    \item \textbf{B) Because it directly influences how far a point must be from the majority to be flagged as an outlier} % CORRECT
    \item C) Because it sets the threshold for the number of clusters the algorithm will create
    \item D) Because it decides whether bitmap indexing can be applied efficiently

\end{enumerate}

\begin{tabular}{@{}l p{10.5cm}@{}}
\textbf{Answer:} & \textbf{B} \hspace{1em} Bloom: Understand \hspace{1em} Comp: 0.522s \hspace{1em} Tokens: 497 \\
\textbf{Empirical:} & \textbf{Easy} \hspace{1em} \textbf{Actual (Human):} \underline{\hspace{5cm}}
\end{tabular}

\vspace{0.2cm}
\rule{\textwidth}{0.4pt}
\vspace{0.2cm}


\section*{Question 64}

\textbf{Q:} How does the text justify the need for “appropriate measurement of deviation” before applying outlier detection?

\begin{enumerate}[label=\textbf{\alph{*})}, nosep]
    \item A) Because it determines whether bitmap indexing can be used
    \item \textbf{B) Because different applications may define “outlier” differently, affecting detection performance} % CORRECT
    \item C) Because it sets the number of clusters the algorithm must generate
    \item D) Because it is required for constructing a correlation matrix among attributes

\end{enumerate}

\begin{tabular}{@{}l p{10.5cm}@{}}
\textbf{Answer:} & \textbf{B} \hspace{1em} Bloom: Understand \hspace{1em} Comp: 0.673s \hspace{1em} Tokens: 586 \\
\textbf{Empirical:} & \textbf{Medium} \hspace{1em} \textbf{Actual (Human):} \underline{\hspace{5cm}}
\end{tabular}

\vspace{0.2cm}
\rule{\textwidth}{0.4pt}
\vspace{0.2cm}


\section*{Question 65}

\textbf{Q:} Which approach is most suitable for detecting outliers in **multivariate data** by converting the problem to a univariate one?

\begin{enumerate}[label=\textbf{\alph{*})}, nosep]
    \item A) Applying a distance‑based proximity method directly to the multivariate space.
    \item B) Using a clustering‑based method to form clusters and flag points outside any cluster.
    \item \textbf{C) Transforming the multivariate data into a single attribute (e.g., via a statistical summary) and then applying a univariate outlier detector.} % CORRECT
    \item D) Training a classification model with outlier labels on the original multivariate attributes.

\end{enumerate}

\begin{tabular}{@{}l p{10.5cm}@{}}
\textbf{Answer:} & \textbf{C} \hspace{1em} Bloom: Understand \hspace{1em} Comp: 0.608s \hspace{1em} Tokens: 568 \\
\textbf{Empirical:} & \textbf{Medium} \hspace{1em} \textbf{Actual (Human):} \underline{\hspace{5cm}}
\end{tabular}

\vspace{0.2cm}
\rule{\textwidth}{0.4pt}
\vspace{0.2cm}


\section*{Question 66}

\textbf{Q:} Why is it important to include a contrasting measure such as count\% when presenting the final class comparison of discovered patterns?

\begin{enumerate}[label=\textbf{\alph{*})}, nosep]
    \item A) It reduces the computational cost of mining by limiting the number of rules.
    \item \textbf{B) It provides a quantitative way to compare the prevalence of a pattern in the target class versus a contrasting class, highlighting its discriminative power.} % CORRECT
    \item C) It eliminates the need for attribute relevance analysis.
    \item D) It ensures that all attributes are treated equally in the final description.

\end{enumerate}

\begin{tabular}{@{}l p{10.5cm}@{}}
\textbf{Answer:} & \textbf{B} \hspace{1em} Bloom: Understand \hspace{1em} Comp: 0.550s \hspace{1em} Tokens: 486 \\
\textbf{Empirical:} & \textbf{Medium} \hspace{1em} \textbf{Actual (Human):} \underline{\hspace{5cm}}
\end{tabular}

\vspace{0.2cm}
\rule{\textwidth}{0.4pt}
\vspace{0.2cm}


\section*{Question 67}

\textbf{Q:} How does the “count\%” contrasting measure assist a user in interpreting the interestingness of a rule?

\begin{enumerate}[label=\textbf{\alph{*})}, nosep]
    \item A) It shows the absolute number of instances covered by the rule in the whole database.
    \item \textbf{B) It indicates the proportion of the target class covered relative to the contrasting class, highlighting the rule’s discriminative strength.} % CORRECT
    \item C) It provides a ranking of attributes by their correlation values.
    \item D) It replaces the need for confidence and support calculations.

\end{enumerate}

\begin{tabular}{@{}l p{10.5cm}@{}}
\textbf{Answer:} & \textbf{B} \hspace{1em} Bloom: Understand \hspace{1em} Comp: 0.622s \hspace{1em} Tokens: 559 \\
\textbf{Empirical:} & \textbf{Medium} \hspace{1em} \textbf{Actual (Human):} \underline{\hspace{5cm}}
\end{tabular}

\vspace{0.2cm}
\rule{\textwidth}{0.4pt}
\vspace{0.2cm}


\section*{Question 68}

\textbf{Q:} Why might a user sample the working relation W when scanning for distinct attribute values during generalization?

\begin{enumerate}[label=\textbf{\alph{*})}, nosep]
    \item A) To ensure that every possible value is captured for accurate pattern evaluation.
    \item \textbf{B) To reduce computational effort while still obtaining a representative set of values for evaluating pattern interestingness.} % CORRECT
    \item C) To eliminate the need for any statistical analysis later in the process.
    \item D) To automatically generate contrasting measures for each attribute.

\end{enumerate}

\begin{tabular}{@{}l p{10.5cm}@{}}
\textbf{Answer:} & \textbf{B} \hspace{1em} Bloom: Understand \hspace{1em} Comp: 0.591s \hspace{1em} Tokens: 541 \\
\textbf{Empirical:} & \textbf{Medium} \hspace{1em} \textbf{Actual (Human):} \underline{\hspace{5cm}}
\end{tabular}

\vspace{0.2cm}
\rule{\textwidth}{0.4pt}
\vspace{0.2cm}
\end{document}